% Options for packages loaded elsewhere
\PassOptionsToPackage{unicode}{hyperref}
\PassOptionsToPackage{hyphens}{url}
\PassOptionsToPackage{dvipsnames,svgnames,x11names}{xcolor}
%
\documentclass[
  11pt,
  a4paper,
]{article}

\usepackage{amsmath,amssymb}
\usepackage{setspace}
\usepackage{iftex}
\ifPDFTeX
  \usepackage[T1]{fontenc}
  \usepackage[utf8]{inputenc}
  \usepackage{textcomp} % provide euro and other symbols
\else % if luatex or xetex
  \usepackage{unicode-math}
  \defaultfontfeatures{Scale=MatchLowercase}
  \defaultfontfeatures[\rmfamily]{Ligatures=TeX,Scale=1}
\fi
\usepackage{lmodern}
\ifPDFTeX\else  
    % xetex/luatex font selection
\fi
% Use upquote if available, for straight quotes in verbatim environments
\IfFileExists{upquote.sty}{\usepackage{upquote}}{}
\IfFileExists{microtype.sty}{% use microtype if available
  \usepackage[]{microtype}
  \UseMicrotypeSet[protrusion]{basicmath} % disable protrusion for tt fonts
}{}
\makeatletter
\@ifundefined{KOMAClassName}{% if non-KOMA class
  \IfFileExists{parskip.sty}{%
    \usepackage{parskip}
  }{% else
    \setlength{\parindent}{0pt}
    \setlength{\parskip}{6pt plus 2pt minus 1pt}}
}{% if KOMA class
  \KOMAoptions{parskip=half}}
\makeatother
\usepackage{xcolor}
\usepackage[top=2.5cm,bottom=2.5cm,left=2.5cm,right=2.5cm]{geometry}
\ifLuaTeX
  \usepackage{luacolor}
  \usepackage[soul]{lua-ul}
\else
  \usepackage{soul}
  
\fi
\setlength{\emergencystretch}{3em} % prevent overfull lines
\setcounter{secnumdepth}{3}


\providecommand{\tightlist}{%
  \setlength{\itemsep}{0pt}\setlength{\parskip}{0pt}}\usepackage{longtable,booktabs,array}
\usepackage{calc} % for calculating minipage widths
% Correct order of tables after \paragraph or \subparagraph
\usepackage{etoolbox}
\makeatletter
\patchcmd\longtable{\par}{\if@noskipsec\mbox{}\fi\par}{}{}
\makeatother
% Allow footnotes in longtable head/foot
\IfFileExists{footnotehyper.sty}{\usepackage{footnotehyper}}{\usepackage{footnote}}
\makesavenoteenv{longtable}
\usepackage{graphicx}
\makeatletter
\def\maxwidth{\ifdim\Gin@nat@width>\linewidth\linewidth\else\Gin@nat@width\fi}
\def\maxheight{\ifdim\Gin@nat@height>\textheight\textheight\else\Gin@nat@height\fi}
\makeatother
% Scale images if necessary, so that they will not overflow the page
% margins by default, and it is still possible to overwrite the defaults
% using explicit options in \includegraphics[width, height, ...]{}
\setkeys{Gin}{width=\maxwidth,height=\maxheight,keepaspectratio}
% Set default figure placement to htbp
\makeatletter
\def\fps@figure{htbp}
\makeatother

\usepackage{booktabs}
\usepackage{longtable}
\usepackage{array}
\usepackage{multirow}
\usepackage{wrapfig}
\usepackage{float}
\usepackage{colortbl}
\usepackage{pdflscape}
\usepackage{tabu}
\usepackage{threeparttable}
\usepackage{threeparttablex}
\usepackage[normalem]{ulem}
\usepackage{makecell}
\usepackage{xcolor}
\usepackage{orcidlink}
\definecolor{mypink}{RGB}{219, 48, 122}
\makeatletter
\@ifpackageloaded{caption}{}{\usepackage{caption}}
\AtBeginDocument{%
\ifdefined\contentsname
  \renewcommand*\contentsname{Table of contents}
\else
  \newcommand\contentsname{Table of contents}
\fi
\ifdefined\listfigurename
  \renewcommand*\listfigurename{List of Figures}
\else
  \newcommand\listfigurename{List of Figures}
\fi
\ifdefined\listtablename
  \renewcommand*\listtablename{List of Tables}
\else
  \newcommand\listtablename{List of Tables}
\fi
\ifdefined\figurename
  \renewcommand*\figurename{Figure}
\else
  \newcommand\figurename{Figure}
\fi
\ifdefined\tablename
  \renewcommand*\tablename{Table}
\else
  \newcommand\tablename{Table}
\fi
}
\@ifpackageloaded{float}{}{\usepackage{float}}
\floatstyle{ruled}
\@ifundefined{c@chapter}{\newfloat{codelisting}{h}{lop}}{\newfloat{codelisting}{h}{lop}[chapter]}
\floatname{codelisting}{Listing}
\newcommand*\listoflistings{\listof{codelisting}{List of Listings}}
\usepackage{amsthm}
\theoremstyle{plain}
\newtheorem{theorem}{Theorem}[section]
\theoremstyle{plain}
\newtheorem{proposition}{Proposition}[section]
\theoremstyle{remark}
\AtBeginDocument{\renewcommand*{\proofname}{Proof}}
\newtheorem*{remark}{Remark}
\newtheorem*{solution}{Solution}
\newtheorem{refremark}{Remark}[section]
\newtheorem{refsolution}{Solution}[section]
\makeatother
\makeatletter
\makeatother
\makeatletter
\@ifpackageloaded{caption}{}{\usepackage{caption}}
\@ifpackageloaded{subcaption}{}{\usepackage{subcaption}}
\makeatother
\ifLuaTeX
  \usepackage{selnolig}  % disable illegal ligatures
\fi
\usepackage[style=authoryear-comp,]{biblatex}
\addbibresource{references.bib}
\usepackage{bookmark}

\IfFileExists{xurl.sty}{\usepackage{xurl}}{} % add URL line breaks if available
\urlstyle{same} % disable monospaced font for URLs
\hypersetup{
  pdftitle={Online conformal inference for multi-step time series forecasting},
  pdfauthor={Xiaoqian Wang; Rob J Hyndman},
  pdfkeywords={Conformal prediction, Coverage
guarantee, Distribution-free inference, Exchangeability, Weighted
quantile estimate},
  colorlinks=true,
  linkcolor={blue},
  filecolor={Maroon},
  citecolor={Blue},
  urlcolor={Blue},
  pdfcreator={LaTeX via pandoc}}

%% CAPTIONS
\usepackage{caption}
\DeclareCaptionStyle{italic}[justification=centering]
 {labelfont={bf},textfont={it},labelsep=colon}
\captionsetup[figure]{style=italic,format=hang,singlelinecheck=true}
\captionsetup[table]{style=italic,format=hang,singlelinecheck=true}

%% FONT
\usepackage{bera}
\usepackage[charter]{mathdesign}
\usepackage[scale=0.9]{sourcecodepro}
\usepackage[lf,t]{FiraSans}

%% HEADERS AND FOOTERS
\usepackage{fancyhdr}
\pagestyle{fancy}
\rfoot{\Large\sffamily\raisebox{-0.1cm}{\textbf{\thepage}}}
\makeatletter
\lhead{\textsf{\expandafter{\@title}}}
\makeatother
\rhead{}
\cfoot{}
\setlength{\headheight}{15pt}
\renewcommand{\headrulewidth}{0.4pt}
\renewcommand{\footrulewidth}{0.4pt}
\fancypagestyle{plain}{%
\fancyhf{} % clear all header and footer fields
\fancyfoot[C]{\sffamily\thepage} % except the center
\renewcommand{\headrulewidth}{0pt}
\renewcommand{\footrulewidth}{0pt}}

%% MATHS
\usepackage{bm,amsmath}
\allowdisplaybreaks

%% GRAPHICS
\makeatletter
\def\fps@figure{htbp}
\makeatother
\setcounter{topnumber}{2}
\setcounter{bottomnumber}{2}
\setcounter{totalnumber}{4}
\renewcommand{\topfraction}{0.85}
\renewcommand{\bottomfraction}{0.85}
\renewcommand{\textfraction}{0.15}
\renewcommand{\floatpagefraction}{0.8}

%% SECTION TITLES
\usepackage[compact,sf,bf]{titlesec}
\titleformat{\section}[block]
  {\fontsize{15}{17}\bfseries\sffamily}
  {\thesection}
  {0.4em}{}
\titleformat{\subsection}[block]
  {\fontsize{12}{14}\bfseries\sffamily}
  {\thesubsection}
  {0.4em}{}
\titlespacing{\section}{0pt}{*5}{*1}
\titlespacing{\subsection}{0pt}{*2}{*0.2}

%% BIBLIOGRAPHY.

\makeatletter
\@ifpackageloaded{biblatex}{
\ExecuteBibliographyOptions{bibencoding=utf8,minnames=1,maxnames=3, maxbibnames=99,dashed=false,terseinits=true,giveninits=true,uniquename=false,uniquelist=false,doi=false, isbn=false,url=true,sortcites=false}
\DeclareFieldFormat{url}{\texttt{\url{#1}}}
\DeclareFieldFormat[article]{pages}{#1}
\DeclareFieldFormat[inproceedings]{pages}{\lowercase{pp.}#1}
\DeclareFieldFormat[incollection]{pages}{\lowercase{pp.}#1}
\DeclareFieldFormat[article]{volume}{\mkbibbold{#1}}
\DeclareFieldFormat[article]{number}{\mkbibparens{#1}}
\DeclareFieldFormat[article]{title}{\MakeCapital{#1}}
\DeclareFieldFormat[article]{url}{}
\DeclareFieldFormat[inproceedings]{title}{#1}
\DeclareFieldFormat{shorthandwidth}{#1}
\usepackage{xpatch}
\xpatchbibmacro{volume+number+eid}{\setunit*{\adddot}}{}{}{}
% Remove In: for an article.
\renewbibmacro{in:}{%
  \ifentrytype{article}{}{%
  \printtext{\bibstring{in}\intitlepunct}}}
\AtEveryBibitem{\clearfield{month}}
\AtEveryCitekey{\clearfield{month}}
\DeclareDelimFormat[cbx@textcite]{nameyeardelim}{\addspace}
\renewcommand*{\finalnamedelim}{\addspace\&\space}
}{}
\makeatother

%% PAGE BREAKING to avoid widows and orphans
\clubpenalty = 2000
\widowpenalty = 2000
\usepackage{microtype}% Placement of logos

\RequirePackage[absolute,overlay]{textpos}
\setlength{\TPHorizModule}{1cm}
\setlength{\TPVertModule}{1cm}
\def\placefig#1#2#3#4{\begin{textblock}{.1}(#1,#2)\rlap{\includegraphics[#3]{#4}}\end{textblock}}

% Title and date

\title{Online conformal inference for multi-step time series
forecasting}
\date{10 September 2024}

\def\Date{\number\day}
\def\Month{\ifcase\month\or
 January\or February\or March\or April\or May\or June\or
 July\or August\or September\or October\or November\or December\fi}
\def\Year{\number\year}

% Working paper number and JEL codes

\makeatletter
\def\wp#1{\gdef\@wp{#1}}\def\@wp{??/??}
\def\jel#1{\gdef\@jel{#1}}\def\@jel{??}
\def\showjel{{\large\textsf{\textbf{JEL classification:}}~\@jel}}
\def\nojel{\def\showjel{}}
\makeatother

\wp{no/yr}
\jel{C10,C14,C32}

% Title page

\makeatletter
\def\cover{{\sffamily\setcounter{page}{0}
        \thispagestyle{empty}
        \placefig{2}{1.5}{width=5cm}{monash2}
        \placefig{16.9}{1.5}{width=2.1cm}{MBSportrait}
        \begin{textblock}{4}(16.9,4)ISSN 1440-771X\end{textblock}
        \begin{textblock}{7}(12.7,27.9)\hfill
        \includegraphics[height=0.7cm]{AACSB}~~~
        \includegraphics[height=0.7cm]{EQUIS}~~~
        \includegraphics[height=0.7cm]{AMBA}
        \end{textblock}
        \vspace*{2.5cm}
        \begin{center}\Large
        Department of Econometrics and Business Statistics\\[.5cm]
        \footnotesize http://monash.edu/business/ebs/research/publications
        \end{center}\vspace{2cm}
        \begin{center}
        \fbox{\parbox{14cm}{\begin{onehalfspace}\centering\Huge\vspace*{0.3cm}
                \textsf{\textbf{\expandafter{\@title}}}\vspace{1cm}\par
                \LARGE
                \expandafter{\@author}
                \end{onehalfspace}
        }}
        \end{center}
        \vfill
                \begin{center}\Large
                \Month~\Year\\[1cm]
                Working Paper \@wp
        \end{center}\vspace*{2cm}}}
        \def\addresses#1{\gdef\@addresses{#1}}\def\@addresses{??}
        \def\pageone{{\sffamily\setstretch{1}%
        \thispagestyle{empty}%
        \vbox to \textheight{%
        \raggedright\baselineskip=1.2cm
     {\fontsize{24.88}{30}\sffamily\textbf{\expandafter{\@title}}}
        \vspace{2cm}\par
        \hspace{1cm}\parbox{14cm}{\sffamily\large\@addresses}\vspace{1cm}\vfill
        \hspace{1cm}{\large\Date~\Month~\Year}\\[1cm]
        \hspace{1cm}\showjel\vss}}}
\def\blindtitle{{\sffamily
     \thispagestyle{plain}\raggedright\baselineskip=1.2cm
     {\fontsize{24.88}{30}\sffamily\textbf{\expandafter{\@title}}}\vspace{1cm}\par
        }}
\def\titlepage{{\cover\newpage\pageone\newpage\blindtitle}}

\def\blind{\def\titlepage{{\blindtitle}}\let\maketitle\blindtitle}
\def\titlepageonly{\def\titlepage{{\pageone\end{document}}}}
\def\nocover{\def\titlepage{{\pageone\newpage\blindtitle}}\let\maketitle\titlepage}
\let\maketitle\titlepage
\makeatother

% Authors


  \author{Xiaoqian Wang, Rob J Hyndman}
  \addresses{%
    %
      \textbf{Xiaoqian Wang}\\%
      %
        Department of Econometrics \& Business Statistics\\%
        Monash University\\%
        Clayton VIC 3800\\%
        Australia\\%
      %
      {Email: xiaoqian.wang@monash.edu}\\%
      \textit{Corresponding author}\\[0.5cm]%
   %
      \textbf{Rob J Hyndman}\\%
      %
        Department of Econometrics \& Business Statistics\\%
        Monash University\\%
        Clayton VIC 3800\\%
        Australia\\%
      %
      {Email: rob.hyndman@monash.edu}\\%
      \\[0.5cm]%
   %
   }%
   \lfoot{\sf Wang, Hyndman: 10 September 2024}

% Keywords

\newenvironment{keywords}{\par\vspace{0.5cm}\noindent{\sffamily\textbf{Keywords:}}}{\vspace{0.25cm}\par\hrule\vspace{0.5cm}\par}

% Abstract
\renewenvironment{abstract}{\begin{minipage}{\textwidth}\parskip=1.4ex\noindent
\hrule\vspace{0.1cm}\par{\sffamily\textbf{\abstractname}}\newline\setstretch{1}}
  {\end{minipage}}
\begin{document}
\maketitle

\begin{abstract}
A brief summary
\end{abstract}

\begin{keywords}
  Conformal prediction; Coverage guarantee; Distribution-free
inference; Exchangeability; 
  Weighted quantile estimate.
\end{keywords}

\setstretch{1}
\section{Introduction}\label{sec-intro}

Motivation: Why conformal prediction?

\begin{itemize}
\tightlist
\item
  Machine Learning, non-parametric models.
\end{itemize}

Background: Development and popularity of conformal prediction

Application to time series data

Challenges and limitations

We consider a general sequential setting in which we observe a time
series \(\{y_t\}_{t \geq 1}\) generated by an unknown data generating
process (DGP), which may depend on its own past, along with other
exogenous predictors, \(\bm{x}_t=(x_{1,t},\ldots,x_{p,t})^{\prime}\),
and their histories. The distribution of
\(\{(\bm{x}_t, y_t)\}_{t \geq 1} \subseteq \mathbb{R}^p \times \mathbb{R}\)
is obviously allowed to vary over time in time series context. At each
time point \(t\), we aim to forecast \(H\) steps into the future,
providing a \emph{prediction set} (which is a prediction interval in
this setting), \(\hat{\mathcal{C}}_{t+h|t}\), for the realization
\(y_{t+h}\) for each \(h\in[H]\). The \(h\)-step-ahead forecast uses the
previously observed data \(\{(\bm{x}_i, y_i)\}_{1 \leq i \leq t}\) along
with the new information of the exogenous predictors
\(\{\bm{x}_{t+j}\}_{1\leq j\leq h}\). Note that we can generate ex-ante
forecasts by using forecasts of the predictors based on information
available up to and including time \(t\). Alternatively, ex-post
forecasts are generated assuming that actual observations of the
predictors from the forecast period are available. Given a nominal
\emph{miscoverage rate} \(\alpha \in (0,1)\) specified by the user, we
expect the output \(\hat{\mathcal{C}}_{t+h|t}\) to be a \emph{valid}
prediction interval so that \(y_{t+h}\) falls within the prediction
interval \(\hat{\mathcal{C}}_{t+h|t}\) at least \(100(1-\alpha)\%\) of
the time.

Contribution: Overview of proposed methods

Some notes such as we use split conformal prediction throughout this
paper, and uniform notation.

Paper structure.

\section{Related work}\label{sec-review}

\subsection{Conformal prediction for regression}\label{sec-review_reg}

In this section, we focus on the regression setting, which stands as one
of the primary areas where conformal prediction methods have seen
substantial development. Suppose we have \(n\) data points
\(Z_i = (X_i, Y_i) \in \mathbb{R}^d \times \mathbb{R}\),
\(i=1,\ldots,n\). The aim of conformal prediction is to construct a
prediction interval
\(\hat{\mathcal{C}}_{n+1}^{\alpha}\left(X_{n+1}\right)\) ensuring that
the unseen response \(Y_{n+1}\) falls within
\(\hat{\mathcal{C}}_{n+1}^{\alpha}\left(X_{n+1}\right)\) at least
\(100(1-\alpha)\%\) of the time.

\subsubsection{Split conformal
prediction}\label{split-conformal-prediction}

\textbf{Split conformal prediction} (SCP, also called inductive
conformal prediction, \textcite{papadopoulos2002}; \textcite{vovk2005};
\textcite{lei2018}), is a holdout method for building prediction
intervals using a pre-trained model.

In regression setting, SCP randomly separates the available \(n\) data
points into a \emph{proper training set} \(\mathcal{D}_{\text{tr}}\) of
size \(n_t\) and a \emph{calibration set} \(\mathcal{D}_{\text{cal}}\)
of size \(n_c\). Given a regression model
\(\hat{\mu}: \mathbb{R}^d \rightarrow \mathbb{R}\) that is fitted on the
training set and a score function \(\mathcal{S}\), a \emph{nonconformity
score} \(s_i = \mathcal{S}\left(X_i, Y_i\right)\),
\(i\in\mathcal{D}_{\text{cal}}\), is computed on every data point in
\(\mathcal{D}_{\text{cal}}\) to measure the nonconformity between the
calibration's response values and the predicted values obtained from the
fitted model \(\hat{\mu}\). One commonly used nonconformity score
function \(\mathcal{S}\) in regression is the absolute residual,
i.e.~\(s_i = |Y_i - \hat{\mu}(X_i)|\). Then SCP computes the prediction
interval for the test data \(Y_{n+1}\) using

\begin{equation}\phantomsection\label{eq-scp}{
\hat{\mathcal{C}}_{n+1}^{\alpha}\left(X_{n+1}\right) = \left\{y\in\mathbb{R}: \mathcal{S}\left(X_{n+1}, y\right) \leq Q_{1-\alpha}\left(\sum_{i \in \mathcal{D}_{\text{cal}}}\frac{1}{n_c+1}\cdot\delta_{s_{i}}+\frac{1}{n_c+1}\cdot\delta_{+\infty}\right)\right\},
}\end{equation}

where \(\mathrm{Q}_\tau(\cdot)\) denotes the \(\tau\)-quantile of its
argument, and \(\delta_a\) denotes the point mass at \(a\).

\begin{theorem}[SCP, \textcite{vovk2005};
\textcite{lei2018}]\protect\hypertarget{thm-scp}{}\label{thm-scp}

Assume that the data points \((X_i, Y_i)\), \(i=1,\ldots,n+1\), are
i.i.d. (or more generally, exchangeable) from any distribution. For any
score function \(\mathcal{S}\), and any \(\alpha\in(0,1)\), the split
conformal prediction interval defined in Equation~\ref{eq-scp} satisfies

\[
\mathbb{P}\left\{Y_{n+1} \in \hat{\mathcal{C}}_{n+1}^{\alpha}\left(X_{n+1}\right)\right\} \geq 1-\alpha.
\]

Moreover, if we assume additionally that the nonconformity scores on
\(\mathcal{D}_{\text{cal}}\) are distinct with probability one, then we
also have

\[
\mathbb{P}\left\{Y_{n+1} \in \hat{\mathcal{C}}_{n+1}^{\alpha}\left(X_{n+1}\right)\right\} < 1-\alpha+\frac{1}{n_c+1}.
\]

\end{theorem}

\subsubsection{Nonexchangeable conformal
prediction}\label{nonexchangeable-conformal-prediction}

\textcite{barber2023} propose \textbf{nonexchangeable conformal
prediction} (NexCP) that generalizes the SCP method to allow for some
sources of nonexchangeability. For split conformal, the NexCP method can
be considered as simply using weighted quantiles to obtain robust
inference. Note that NexCP assumes the weights are fixed and
data-independent. The intuition is that a higher weight should be
assigned to a data point that is believed to originate from the same
distribution as the test data.

Given weights \(w_i \in [0,1]\), \(i \in \mathcal{D}_{\text{cal}}\), the
prediction interval of the NexCP method for the test data \(Y_{n+1}\) is
given by

\begin{equation}\phantomsection\label{eq-nexcp}{
\hat{\mathcal{C}}_{n+1}^{\alpha}\left(X_{n+1}\right) = \left\{y\in\mathbb{R}: \mathcal{S}\left(X_{n+1}, y\right) \leq Q_{1-\alpha}\left(\sum_{i \in \mathcal{D}_{\text{cal}}}\tilde{w}_i\cdot\delta_{s_{i}}+\tilde{w}_{n+1}\cdot\delta_{+\infty}\right)\right\},
}\end{equation}

where \(\tilde{w}_i\) and \(\tilde{w}_{n+1}\) are normalized weights
calculated by

\[
\tilde{w}_i = \frac{w_i}{\sum_{i\in\mathcal{D}_{\text{cal}}}w_i+1}, \text{ for } i \in \mathcal{D}_{\text{cal}} \quad \text{and} \quad \tilde{w}_{n+1} =  \frac{1}{\sum_{i\in\mathcal{D}_{\text{cal}}}w_i+1}.
\]

By placing different prespecified weights on data points, NexCP is able
to deal with data that are not exchangeable and provide robustness
against distribution drift. \textcite{barber2023} suggests using an
exponential weighting scheme for time series data, where the weights
decrease exponentially as data points come from the further in the past.

\begin{theorem}[NexCP,
\textcite{barber2023}]\protect\hypertarget{thm-nexcp}{}\label{thm-nexcp}

Let \(\mathcal{S}(Z)\) denote a vector of nonconformity scores for data
points in the calibration and test sets, and \(\mathcal{S}(Z^i)\) denote
a vector of nonconformity scores after swapping the test point with the
\(i\)th data point in the calibration set.

For any score function \(\mathcal{S}\), and any \(\alpha\in(0,1)\), the
nonexchangeable split conformal prediction interval defined in
Equation~\ref{eq-nexcp} satisfies

\[
\mathbb{P}\left\{Y_{n+1} \in \hat{\mathcal{C}}_{n+1}^{\alpha}\left(X_{n+1}\right)\right\} \geq 1-\alpha-\sum_{i \in \mathcal{D}_{\text{cal}}} \tilde{w}_i \cdot \mathrm{d}_{\mathrm{TV}}\left(\mathcal{S}(Z), \mathcal{S}\left(Z^i\right)\right),
\]

without the assumption of exchangeability of the data, where
\(\mathrm{d}_{\mathrm{TV}}\) denotes the total variation distance
between two distributions.

Moreover, if we assume additionally that the nonconformity scores on
\(\mathcal{D}_{\text{cal}}\) are distinct with probability one, then the
probability also has the upper bound:

\[
\mathbb{P}\left\{Y_{n+1} \in \hat{\mathcal{C}}_{n+1}^{\alpha}\left(X_{n+1}\right)\right\} < 1-\alpha+\tilde{w}_{n+1}+\sum_{i \in \mathcal{D}_{\text{cal}}} \tilde{w}_i \cdot \mathrm{d}_{\mathrm{TV}}\left(\mathcal{S}(Z), \mathcal{S}\left(Z^i\right)\right).
\]

\end{theorem}

\subsubsection{Adaptive conformal
prediction}\label{adaptive-conformal-prediction}

The \textbf{adaptive conformal prediction} (ACP) proposed by
\textcite{gibbs2021} accounts for nonexchangeability by updating the
quantile level in an online manner. Specifically, it treats \(\alpha\)
as a tunable parameter and estimates it recursively based on the
historical performance. ACP can be used to deal with arbitrary online
distribution shifts.

Similar to SCP, the initial step involves a random split on the observed
data, yielding a training set for fitting a regression model and a
withheld calibration set. ACP assumes that there exists an optimal value
\(\alpha_{t}^{*}\) to achieve the desired miscoverage rate \(\alpha\) at
each time \(t\). To deal with cases where the data generating
distribution is shifting over time, ACP recursively estimates the
parameter \(\alpha_t^{*}\) on the test points, using the updating
equation

\begin{equation}\phantomsection\label{eq-acp}{
\alpha_{t+1} = \alpha_{t} + \gamma\left(\alpha - \mathrm{err}_t\right),
}\end{equation}

rather than consistently using the target miscoverage rate \(\alpha\).
Here, \(\gamma > 0\) is a fixed step size parameter, \(\alpha_1\) is the
initial estimate typically set as \(\alpha_1=\alpha\), and
\(\mathrm{err}_t = \mathbb{1}\left\{Y_t \notin \hat{\mathcal{C}}_{t}^{\alpha_t}\left(X_t\right)\right\}\),
where \(\hat{\mathcal{C}}_{t}^{\alpha_t}\left(X_t\right)\) denotes the
prediction set obtained using the \(1-\alpha_{t}\) quantile for the
nonconformity scores available up to and including time \(t\).

This update process adapts the estimation of \(\alpha_{t}^{*}\) based on
the historical frequency of miscoverage in the prediction sets.
Specifically, it adjusts upwards (or downwards) the estimate of
\(\alpha_{t}^{*}\) if the prediction sets have shown over-coverage (or
under-coverage) of the actual outcomes.

Selecting the parameter \(\gamma\) is pivotal yet challenging.
\textcite{gibbs2021} suggests setting \(\gamma\) in proportion to the
degree of variation of the unknown \(\alpha_{t}^{*}\) over time. Several
strategies have been proposed to avoid the necessity of selecting
\(\gamma\). For example, \textcite{zaffran2022} use an adaptive
aggregation of multiple ACPs with a set of candidate values for
\(\gamma\) , determining weights based on their historical performance.
\textcite{bastani2022} propose the multivalid prediction algorithm in
which the prediction set is established by selecting a threshold from a
sequence of candidate thresholds. However, \textcite{gibbs2024} have
empirically shown that both previous methods fail to promptly adapt to
the local changes. To address this limitation, \textcite{gibbs2024}
propose adaptively tuning the step size parameter \(\gamma\) in an
online setting, choosing an ``optimal'' value for \(\gamma\) from a
candidate set of values by assessing their historical performance.

\begin{theorem}[ACP,
\textcite{gibbs2021}]\protect\hypertarget{thm-acp}{}\label{thm-acp}

Assume that, with probability one, \(Q_{\tau, t}\) is continuous and
nondecreasing so that \(Q_{0, t}=-\infty\) and \(Q_{1, t}=\infty\). Then
for any \(\alpha\in(0,1)\), \(\gamma > 0\), and any \(T \geq 1\) test
points, the prediction sets given by ACP satisfy

\[
|\frac{1}{T}\sum_{t=1}^{T}\mathrm{err}_t - \alpha| \leq \frac{\max\left\{\alpha_1,1-\alpha_1\right\}+\gamma}{\gamma T}.
\]

In particular, this means that the prediction intervals obtained by ACP
yield long-run coverage,
i.e.~\(\lim _{T \rightarrow \infty} \frac{1}{T} \sum_{t=1}^T \mathrm{err}_t = \alpha\).

\end{theorem}

\textbf{Remark.} Theorem~\ref{thm-acp} suggests that a larger value for
\(\gamma\) generally results in less deviation from the target coverage.
As there is no restriction on \(\alpha_t\) and it can drift below \(0\)
or above \(1\), a larger \(\gamma\) may lead to frequent output of null
or infinite prediction sets in order to quickly adapt to the current
miscoverage status.

\subsubsection{Conformal PID control}\label{conformal-pid-control}

Instead of iteratively updating the miscoverage rate \(\alpha\) as in
\textcite{gibbs2021}, \textcite{angelopoulos2024} draw inspiration from
control theory and directly update the quantile estimate \(q_t\) in an
online fashion to achieve long-run coverage. This method treats the
system for producing prediction sets as a
\textbf{proportional-integral-derivative} controller, later referred to
as PID.

The iteration of the PID method is given by

\[
q_{t+1}=\underbrace{q_t+\eta \left(\mathrm{err}_t-\alpha\right)}_{\mathrm{P}}+\underbrace{r_t\left(\sum_{i=1}^t \left(\mathrm{err}_i-\alpha\right)\right)}_{\mathrm{I}}+\underbrace{g_t^{\prime}}_{\mathrm{D}}.
\]

The PID method integrates three modules, namely quantile tracking (P
control), error integration (I control), and scorecasting (D control).

The P control module updates the quantile iteratively with a constant
learning rate \(\eta > 0\). The underlying intuition is similar to that
of ACP: it increases (or decreases) the quantile estimate if the
prediction set at time \(t\) miscovered (or covered) the corresponding
realization. ACP can be considered as a special case of the P control,
while the P control has the ability to prevent the generation of null or
infinite prediction sets after a sequence of miscoverage events.

\begin{theorem}[The P control,
\textcite{angelopoulos2024}]\protect\hypertarget{thm-pid_p}{}\label{thm-pid_p}

Assume that the nonconformity scores are bounded within \([-b, b]\), for
\(0<b<\infty\). Then for any \(\alpha \in (0,1)\), \(\eta > 0\), and any
\(T \geq 1\) the P control iteration satisfies

\[
|\frac{1}{T}\sum_{t=1}^{T}\mathrm{err}_t - \alpha| \leq \frac{b+\eta}{\eta T}.
\]

In particular, this means that the prediction intervals obtained by the
P control iteration yield long-run coverage,
i.e.~\(\lim _{T \rightarrow \infty} \frac{1}{T} \sum_{t=1}^T \mathrm{err}_t = \alpha\).

\end{theorem}

The I control involves the sum of historical coverage errors
\(\sum_{i=1}^t (\mathrm{err}_i-\alpha)\) in a saturation function
\(r_t\) when updating the quantile estimate, further stabilizing the
coverage.

\begin{theorem}[The I control,
\textcite{angelopoulos2024}]\protect\hypertarget{thm-pid_i}{}\label{thm-pid_i}

Assume that the nonconformity scores are bounded within \([-b, b]\), for
\(b>0\), and that the saturation function \(r_t\) satisfies

\begin{equation}\phantomsection\label{eq-saturation}{
x \geq c \cdot g(t) \Longrightarrow r_t(x) \geq b, \quad \text {and} \quad x \leq-c \cdot g(t) \Longrightarrow r_t(x) \leq -b,
}\end{equation}

for positive constants \(b, c\) and an admissible function \(g\). Then
for any \(\alpha \in (0,1)\), and any \(T \geq 1\) the I control
iteration satisfies

\[
|\frac{1}{T}\sum_{t=1}^{T}\mathrm{err}_t - \alpha| \leq \frac{c \cdot g(T)+1}{T}.
\]

In particular, this means that the prediction intervals obtained by the
I control iteration yield long-run coverage,
i.e.~\(\lim _{T \rightarrow \infty} \frac{1}{T} \sum_{t=1}^T \mathrm{err}_t = \alpha\).

\end{theorem}

Finally, the D control \(g_t^{\prime}\) is the forecast of \(q_{t+1}\)
produced by a scorecaster (forecasting model) fitted using the
nonconformity scores available up to and including time \(t\). Instead
of reacting to the past miscoverage events, this module looks forward
and identifies the leftover signal not captured by the regression model
\(\hat{\mu}\). However, it may introduce variability and result in wider
prediction sets if the scorecaster is aggressive or if there is not much
leftover signal in the nonconformity scores.

\subsection{Conformal prediction for time series}\label{sec-review_ts}

\begin{itemize}
\item
  Brief literature review of conformal prediction methods or
  applications on time series data.
\item
  Multi-step conformal prediction.
\end{itemize}

\section{Multi-step conformal prediction for time
series}\label{sec-method}

We now consider multi-step time series forecasting problems. In the
following sections, we first introduce the setup for time series
forecasting problems, and then generalize the existing conformal
prediction methods described in Section~\ref{sec-review_reg} to deal
with multi-step time series forecasting. Finally, we explore the
properties of forecast errors for optimal multi-step forecasts, and then
propose the \textbf{autocorrelated multi-step conformal prediction}
(AcMCP) method to account for the serial correlation of multi-step
forecast errors.

\subsection{Setup}\label{sec-setup}

Let \(z_t = (\bm{x}_t, y_t)\) denote the data point (including the
response and possibly predictors) at time \(t\). Suppose that, at each
time \(t\), we have a forecasting model \(\hat{f}_t\) trained using the
historical data \(z_{1:t}\). Throughout the paper, we assume that the
predictors are known into the future. In this way, we perform ex-post
forecasting and there is no additional uncertainty introduced from
forecasting the exogenous predictors. Using the forecasting model
\(\hat{f}_t\), we are able to produce \(H\)-step point forecasts,
\(\{\hat{y}_{t+h|t}\}_{h\in[H]}\), using the future values for the
predictors. The task is to employ conformal inference to build
\(H\)-step prediction intervals,
\(\{\hat{\mathcal{C}}_{t+h|t}^{\alpha}\left(z_{1:t},\bm{x}_{t+1:h}\right)\}_{h\in[H]}\),
at the target coverage level \(1-\alpha\). For brevity, we will use
\(\hat{\mathcal{C}}_{t+h|t}^{\alpha}\) to denote the \(h\)-step-ahead
prediction interval produced using the information up to time \(t\).

\textbf{Sequential split.} In time series context, it is infeasible to
perform random splitting in split conformal due to the temporal
dependency present in the data. Instead, throughout the conformal
prediction methods in this section, we use a sequential split to
preserve the temporal structure. For example, the \(t\) available data
points, \(z_{1:t}\), are sequentially split into two consecutive sets, a
proper training set \(\mathcal{D}_{\text{tr}} \subset \{1,\ldots,t_r\}\)
and a calibration set
\(\mathcal{D}_{\text{cal}} \subset \{t_r+1,\ldots,t\}\), where
\(t-t_r \gg h\).

\textbf{Online learning.} Here we consider a generic online learning
framework to adapt to all conformal prediction methods we will discuss
in subsequent sections. The framework adopts a standard rolling window
evaluation strategy. Let the length of the training set be denoted as
\(t_r\)\hspace{0pt} and the calibration set length as
\(t_c\)\hspace{0pt}. The online learning framework consists of the
following steps.

\begin{enumerate}
\def\labelenumi{\arabic{enumi}.}
\item
  Initialization. Train a forecasting model on the initial proper
  training set \(z_{t-t_r+1:t}\), setting \(t=t_r\). Then generate
  \(H\)-step point forecasts \(\{\hat{y}_{t_r+h|t_r}\}_{h\in[H]}\) and
  compute the corresponding nonconformity scores
  \(\{s_{t_r+h|t_r}=\mathcal{S}(z_{1:t_r}, y_{t_r+h})\}_{h\in[H]}\)
  based on the true values \(H\) time steps ahead,
  i.e.~\(\{y_{t_r+h}\}_{h\in[H]}\).
\item
  Recurring procedure. Roll the training set forward by one data point
  by setting \(t \rightarrow t+1\). Then repeat the step 1 until the
  nonconformity scores on the entire initial calibration set,
  \(\{s_{t+h|t}\}_{t_r \leq t \leq t_r+t_c-h}\) for \(h\in[H]\), are
  computed.
\item
  Quantile estimation and prediction interval calculation. Use
  nonconformity scores obtained from the calibration set to perform
  quantile estimation and compute \(H\)-step prediction intervals on the
  test set.
\item
  Online updating. Recursively roll the training set and calibration set
  forward by one data point to update the nonconformity scores for
  calibration, and then repeat the step 3 until prediction intervals for
  the entire test set are obtained, i.e.,
  \(\{\hat{\mathcal{C}}_{t+h|t}^{\alpha}\}_{t_r+t_c \leq t \leq T-h}\)
  for \(h \in [H]\), where \(T-t_r-t_c\) is the length of the test set.
  Therefore, our goal is to achieve long-run coverage in time.
\end{enumerate}

For simplicity, so far we have only presented the nonconformity score
defined as the (signed) forecast error

\[
s_{t+h|t}=\mathcal{S}\left(z_{1:t}, y_{t+h}\right):=y_{t+h}-\hat{f}_t\left(z_{1:t},\bm{x}_{t+1:h}\right)=y_{t+h}-\hat{y}_{t+h|t},
\]

which is the most commonly used accuracy measure in the context of time
series forecasting. We also note that the online learning setting can
also be easily adjusted to work with expanding windows for the training
and calibration sets.

\textbf{Remark.} With sequential splitting, multiple \(H\)-step
forecasts and their respective nonconformity scores can be computed on
the calibration set. These nonconformity scores have diverse forecast
horizons, ranging from \(1\) to \(H\), i.e., the number of periods
between the forecast origin and the time at which nonconformity scores
are evaluated. Thus, we can not uniformly treat these nonconformity
scores and generate \(H\)-step prediction intervals of identical width.

\subsection{Related methods extensions to multi-step
forecasting}\label{sec-ext}

One of the key properties of optimal forecast errors is that the
variance of the forecast error \(e_{t+h|t}\) is non-decreasing in \(h\)
\autocite{diebold1996,patton2007}. Therefore, instead of uniformly
treating \(H\)-step nonconformity scores and generating \(H\)-step
prediction intervals of identical width, we consider a setting wherein a
separate conformal prediction procedure is applied for each
\(h \in [H]\) in an online manner.

\subsubsection{Online multi-step split conformal
prediction}\label{online-multi-step-split-conformal-prediction}

We introduce online \textbf{multi-step split conformal prediction}
(MSCP) as a generalization of SCP to recursively update all \(H\)-step
prediction intervals over time. Specifically, for each \(h \in [H]\), we
consider the following simple online update to construct prediction
intervals on the test set:

\begin{equation}\phantomsection\label{eq-mscp}{
\hat{\mathcal{C}}_{t+h|t}^{\alpha} = \left\{y\in\mathbb{R}: s_{t+h|t}^{y} \leq Q_{1-\alpha}\left(\sum_{i=t-t_c+1}^{t}\frac{1}{t_c+1}\cdot\delta_{s_{i|i-h}}+\frac{1}{t_c+1}\cdot\delta_{+\infty}\right)\right\},
}\end{equation}

where \(s_{t+h|t}^{y}:=\mathcal{S}(z_{1:t}, y)\) denote the
\(h\)-step-ahead nonconformity score calculated at time \(t\) using a
hypothesized test observation \(y\).

\begin{itemize}
\item
  Quantile estimation.
\item
  Theorem?
\end{itemize}

\subsubsection{Online multi-step weighted conformal
prediction}\label{online-multi-step-weighted-conformal-prediction}

The online \textbf{multi-step weighted conformal prediction} (MWCP)
method adapts the NexCP method to the online setting. MWCP uses weighted
quantile estimate for constructing prediction intervals, contrasting
with the MSCP definitions where all nonconformity scores for calibration
are implicitly assigned equal weight.

We choose fixed weights \(w_i = b^{t+1-i}\), \(b \in (0, 1)\) and
\(i=t-t_c+1,\ldots,t\), for nonconformity scores on the corresponding
calibration set. In this setting, weights decay exponentially as the
nonconformity scores get order, akin to the rationale behind the
exponential smoothing method in time series forecasting. Then for each
\(h \in [H]\), MWCP consider the online update for \(h\)-step-ahead
prediction interval:

\[
\hat{\mathcal{C}}_{t+h|t}^{\alpha} = \left\{y\in\mathbb{R}: s_{t+h|t}^{y} \leq Q_{1-\alpha}\left(\sum_{i=t-t_c+1}^{t}\tilde{w}_i\cdot\delta_{s_{i|i-h}}+\tilde{w}_{t+1}\cdot\delta_{+\infty}\right)\right\},
\]

where \(\tilde{w}_i\) and \(\tilde{w}_{t+1}\) are normalized weights
given by

\[
\tilde{w}_i = \frac{w_i}{\sum_{i=t-t_c+1}^{t}w_i+1}, \text{ for } i \in \{t-t_c+1,\ldots,t\} \quad \text{and} \quad \tilde{w}_{t+1} =  \frac{1}{\sum_{i=t-t_c+1}^{t}w_i+1}.
\]

\begin{itemize}
\item
  Weighted quantile estimation.
\item
  Theorem?
\end{itemize}

\subsubsection{Multi-step adaptive conformal
prediction}\label{multi-step-adaptive-conformal-prediction}

In the online learning framework outlined in Section~\ref{sec-setup}, we
extend the ACP method to address multi-step time series forecasting,
introducing the \textbf{multi-step adaptive conformal prediction} (MACP)
method. Specifically, for each \(h \in [H]\), we iteratively estimate
\(\alpha^{*}\) (treated as a tunable parameter) using the update
equation

\begin{equation}\phantomsection\label{eq-macp}{
\alpha_{t+h|t} := \alpha_{t+h-1|t-1} + \gamma\left(\alpha - \mathrm{err}_{t|t-h}\right),
}\end{equation}

and compute the \(h\)-step-ahead prediction interval using
Equation~\ref{eq-mscp} by setting \(\alpha = \alpha_{t+h|t}\). Here,
\(\gamma > 0\) denotes a fixed step size parameter, \(\alpha_{h+1|1}\)
denotes the initial estimate typically set to \(\alpha\), and
\(\mathrm{err}_{t|t-h}\) denotes the miscoverage event
\(\mathrm{err}_{t|t-h} = \mathbb{1}\left\{y_t \notin \hat{\mathcal{C}}_{t|t-h}^{\alpha_{t|t-h}}\right\}\).

Equation~\ref{eq-mscp} indicates that the correction to the estimation
of \(\alpha\) at time \(t+h\) is determined by the historical
miscoverage frequency up to time \(t\). At each iteration, we raise the
estimate of \(\alpha\) used for quantile estimation at time \(t+h\) if
\(\hat{\mathcal{C}}_{t|t-h}^{\alpha_{t|t-h}}\) covered \(y_t\), whereas
we lower the estimate if \(\hat{\mathcal{C}}_{t|t-h}^{\alpha_{t|t-h}}\)
miscovered \(y_t\). Thus the miscoverage event has a delayed impact on
the estimation of \(\alpha\) over \(h\) periods, indicating that the
correction of the \(\alpha\) estimate becomes less prompt with
increasing values of \(h\). Particularly, Equation~\ref{eq-macp} reduces
to the update for ACP as given by Equation~\ref{eq-acp} for \(h=1\).

It should be noted that the update equation
\(\alpha_{t+1|t-h+1} := \alpha_{t|t-h} + \gamma\left(\alpha - \mathrm{err}_{t|t-h}\right)\)
is not utilized in this context. This arises from its limitation that
the available information at time \(t\) is insufficient to derive the
\(\alpha\) estimate for forecasting \(h\) steps ahead.

\begin{itemize}
\tightlist
\item
  Theorem?
\end{itemize}

\subsubsection{Multi-step conformal PID
control}\label{multi-step-conformal-pid-control}

We propose \textbf{multi-step conformal PID control} method (referred to
as MPID hereafter), a generalization of the PID method to deal with
multi-step forecasting.

For each individual forecast horizon \(h\in[H]\), the iteration of the
\(h\)-step-ahead quantile estimate is given by

\[
q_{t+h|t}=\underbrace{q_{t+h-1|t-1}+\eta \left(\mathrm{err}_{t|t-h}-\alpha\right)}_{\mathrm{P}}+\underbrace{r_t\left(\sum_{i=h+1}^t \left(\mathrm{err}_{i|i-h}-\alpha\right)\right)}_{\mathrm{I}}+\underbrace{\hat{s}_{t+h|t}}_{\mathrm{D}},
\]

where as before, \(\eta > 0\) is a constant learning rate, and \(r_t\)
is a saturation function that adheres to the following conditions

\begin{equation}\phantomsection\label{eq-saturation_h}{
x \geq c \cdot g(t-h) \Longrightarrow r_t(x) \geq b, \quad \text {and} \quad x \leq-c \cdot g(t-h) \Longrightarrow r_t(x) \leq -b,
}\end{equation}

for constant \(b, c > 0\), and an admissible function \(g\) that is
sublinear, nonnegative, nondecreasing. Here, we define
\(\hat{s}_{t+h|t}\) as the \(h\)-step-ahead forecast of the
nonconformity score (defined as the forecast error), produced by any
suitable scorecaster (forecasting model) trained using the
\(h\)-step-ahead nonconformity scores available up to and including time
\(t\). With this updating equation, we can obtain all required
\(h\)-step-ahead prediction intervals using information available up to
time \(t\).

The P control in MPID shows a delayed correction of the \(\alpha\)
estimate for a length of \(h\) periods. The I control accounts for the
cumulative historical coverage errors associated with \(h\)-step-ahead
prediction intervals during the update process, thereby enhancing the
stability of the interval coverage. Moreover, the D control performs
\(h\)-step-ahead forecasting, which tends to result in increased
forecast variance for larger forecast horizon \(h\).

\begin{itemize}
\tightlist
\item
  Theorem?
\end{itemize}

\subsection{Autocorrelated multi-step conformal
prediction}\label{sec-acmcp}

In the PID method proposed by \textcite{angelopoulos2024}, a notable
feature is the inclusion of a scorecaster, a model trained on the score
sequence, to forecast the next score. The rationale behind it is to
residualize out any leftover signal in the score distribution not
captured by the base forecasting model, such as trend and seasonality.
However, in the context of time series forecasting, good forecasts are
essential for making good decisions. We naturally expect to use a good
forecasting model and ensure there is no useful signal in forecast
errors (defined as nonconformity scores in this paper). If the forecasts
are not optimal, the forecasting model should be improved to enhance its
performance. Hence, we typically assume the use of a good forecasting
model, and therefore, relying on another model to predict forecast
errors to capture leftover information is not a commonly used solution.
Moreover, the inclusion of a scorecaster often only introduces variance
to the quantile estimate, resulting in wider prediction intervals.

On the other hand, in our general setup outlined in
Section~\ref{sec-intro} and Section~\ref{sec-setup}, the DGP of a time
series may depend on its own past, along with other exogenous predictors
and their histories. Consequently, the \(h\)-step-ahead forecast errors
\(e_{t+h|t}\) may depend on the forecast errors from the past \(h-1\)
steps, i.e.~\(e_{t+1|t}, \ldots, e_{t+h-1|t}\), and forecast errors may
accumulate over the forecast horizon. However, no conformal prediction
methods have taken this potential dependence into account in their
methodological construction.

In this section, we will explore the properties of multi-step forecast
errors and propose a novel conformal prediction method that considers
the autocorrelations of multi-step forecast errors.

\subsubsection{Properties of multi-step forecast errors}\label{sec-ppt}

We assume that a time series \(\{y_t\}_{t \geq 1}\) is generated by a
general non-stationary autoregressive process given by:

\begin{equation}\phantomsection\label{eq-dgp}{
y_t = f_{t}\left(y_{(t-d):(t-1)},\bm{x}_{(t-k):t}\right) + \epsilon_t,
}\end{equation}

where \(f_{t}\) is considered a nonlinear function in \(d\) lagged
values of \(y_t\) (i.e.~\(y_{(t-d):(t-1)}\)) and the current value along
with the preceding \(k\) values of the exogenous predictors
(i.e.~\(\bm{x}_{(t-k):t}\)), and \(\epsilon_t\) is white noise.

It is well-established in the forecasting literature that, for optimal
\(h\)-step-ahead forecasts, the sequence of forecast errors is \emph{at
most} an MA\((h-1)\) process \autocite{harvey1997,diebold2017}. We now
present the property under the assumption of a non-stationary
autoregressive DGP, and provide its proof in Section~\ref{sec-proof_ma}
based on the proof of Proposition~\ref{prp-ar} that we will introduce
later.

\begin{proposition}[MA\((h-1)\) process for \(h\)-step-ahead optimal
forecast errors]\protect\hypertarget{prp-ma}{}\label{prp-ma}

Let \(\{y_t\}_{t \geq 1}\) be a time series generated by a general
non-stationary autoregressive process as given in Equation~\ref{eq-dgp}.
Assume that the exogenous predictors are known into the future if
applicable. The forecast errors of optimal \(h\)-step-ahead forecasts
follow an approximate MA(\(h-1\)) process:

\[
e_{t+h|t} = c + \epsilon_{t+h} + \theta_1\epsilon_{t+h-1} + \cdots + \theta_{h-1}\epsilon_{t+1},
\]

where \(c=0\), motivated by the property that optimal forecasts are
unbiased.

\end{proposition}

We proceed by exploring the autocorrelations of multi-step forecast
errors for optimal forecasts.

\begin{proposition}[Autocorrelations of multi-step optimal forecast
errors]\protect\hypertarget{prp-ar}{}\label{prp-ar}

Let \(\{y_t\}_{t \geq 1}\) be a time series generated by a general
non-stationary autoregressive process as given in Equation~\ref{eq-dgp}.
Assume that the exogenous predictors are known into the future if
applicable. The forecast errors of optimal \(h\)-step-ahead forecasts
are \ul{at most} an approximate AR(\(h-1\)) process given by:

\[
e_{t+h|t} = c + \epsilon_{t+h} + \phi_1e_{t+h-1|t} + \cdots + \phi_{h-1}e_{t+1|t},
\]

where \(e_{t+h|t}\) is the \(h\)-step-ahead forecast error with variance
non-decreasing in \(h\), and the intercept \(c=0\), given the property
that optimal forecasts are unbiased.

\end{proposition}

\begin{proof}
Let \(\hat{y}_{t+h|t}\) be the optimal \(h\)-step-ahead point forecast,
and \(e_{t+h|t}\) be the \(h\)-step-ahead forecast error. Denote
\(\bm{u}_{t+h}=\bm{x}_{(t-k+h):(t+h)}\). Then we have

\[
\begin{aligned}
\hat{y}_{t+h|t}=\begin{cases}
      f_{t+1}\left(y_t,\cdots,y_{t-d+1},\bm{u}_{t+1}\right) & \text{ if } h=1, \\
      f_{t+h}\left(\hat{y}_{t+h-1|t},\cdots,\hat{y}_{t+1|t},y_t,\cdots,y_{t+h-d},\bm{u}_{t+h}\right) &  \text{ if } 1 < h \leq d, \\
      f_{t+h}\left(\hat{y}_{t+h-1|t},\cdots,\hat{y}_{t+h-d|t},\bm{u}_{t+h}\right) & \text{ if } h > d.\\
    \end{cases}\\
\end{aligned}
\]

For \(h=1\), we simply have \(e_{t+1|t} = \epsilon_{t+1}\).

For \(1<h\leq d\), applying the first order Taylor series expansion, we
can write

\[
\begin{aligned}
y_{t+h}
=&f_{t+h}\left(y_{t+h-1},\cdots,y_{t+h-d},\bm{u}_{t+h}\right)+\epsilon_{t+h} \\
=&f_{t+h}\left(\hat{y}_{t+h-1|t}+e_{t+h-1|t},\cdots,\hat{y}_{t+1|t}+e_{t+1|t},y_{t},\cdots,y_{t+h-d},\bm{u}_{t+h}\right)+\epsilon_{t+h} \\
\underset{\text{te}}{\approx}&f_{t+h}\left(\bm{a}\right)+\operatorname{D}f_{t+h}\left(\bm{a}\right)\left(\bm{v}-\bm{a}\right)+
\epsilon_{t+h} \\
=&f_{t+h}\left(\hat{y}_{t+h-1|t},\cdots,\hat{y}_{t+1|t},y_{t},\cdots,y_{t+h-d},\bm{u}_{t+h}\right) \\
&+e_{t+h-1|t}\frac{\partial f_{t+h}\left(\bm{a}\right)}{\partial v_1}+\cdots+e_{t+2|t}\frac{\partial f_{t+h}\left(\bm{a}\right)}{\partial v_{h-2}}+e_{t+1|t}\frac{\partial f_{t+h}\left(\bm{a}\right)}{\partial v_{h-1}}+\epsilon_{t+h} \\
=&\hat{y}_{t+h|t}+e_{t+h|t},
\end{aligned}
\]

where \(\bm{v}=\left(y_{t+h-1},\cdots,y_{t+h-d},\bm{u}_{t+h}\right)\),
\(\bm{a} =\left(\hat{y}_{t+h-1|t},\cdots,\hat{y}_{t+1|t},y_{t},\cdots,y_{t+h-d},\bm{u}_{t+h}\right)\),
\(\operatorname{D}f_{t+h}\left(\bm{a}\right)\) denotes the matrix of
partial derivative of \(f_{t+h}(\bm{v})\) at \(\bm{v}=\bm{a}\), and
\(\frac{\partial}{\partial v_i}\) denotes the partial derivative with
respect to the \(i\)th component in \(f_{t+h}\).

Similarly, for \(h > d\), we can write

\[
\begin{aligned}
y_{t+h}
=&f_{t+h}\left(y_{t+h-1},\cdots,y_{t+h-d},\bm{u}_{t+h}\right)+\epsilon_{t+h} \\
=&f_{t+h}\left(\hat{y}_{t+h-1|t}+e_{t+h-1|t},\cdots,\hat{y}_{t+h-d|t}+e_{t+h-d|t},\bm{u}_{t+h}\right)+\epsilon_{t+h} \\
\underset{\text{te}}{\approx}&f_{t+h}\left(\bm{a}\right)+\operatorname{D}f_{t+h}\left(\bm{a}\right)\left(\bm{v}-\bm{a}\right)+
\epsilon_{t+h} \\
=&f_{t+h}\left(\hat{y}_{t+h-1|t},\cdots,\hat{y}_{t+h-d|t},\bm{u}_{t+h}\right) \\
&+e_{t+h-1|t}\frac{\partial f_{t+h}\left(\bm{a}\right)}{\partial v_1}+e_{t+h-d|t}\frac{\partial f_{t+h}\left(\bm{a}\right)}{\partial v_{d}}+\epsilon_{t+h} \\
=&\hat{y}_{t+h|t}+e_{t+h|t},
\end{aligned}
\]

Therefore, the forecast errors of optimal \(h\)-step-ahead forecasts are
at most an approximate AR(\(h-1\)) process.
\end{proof}

Proposition~\ref{prp-ar} can be viewed as an extension of
Proposition~\ref{prp-ma}. It suggests that the \(h\)-step ahead forecast
error, \(e_{t+h|t}\), is serially correlated with the forecast errors
from at most the past \(h-1\) steps, i.e.,
\(e_{t+1|t}, \ldots, e_{t+h-1|t}\). However, we note that the
autocorrelation among errors associated with optimal forecasts can not
be used to improve forecasting performance, as it does not incorporate
any information available when the forecast was made. It is reasonable
because if we could forecast the forecast error, we could improve the
forecast, indicating that the initial forecast couldn't have been
optimal.

The proof of Proposition~\ref{prp-ar} suggests that, if \(f_t\) is a
linear autoregressive model, then the AR coefficients are the linear
coefficients of the optimal forecasting model. However, when \(f_t\)
takes on a more complex nonlinear structure, the AR coefficients become
complicated functions of observed data and unobserved model
coefficients.

\subsubsection{The AcMCP method}\label{sec-novel}

Inspired by the properties of multi-step forecast errors discussed in
Section~\ref{sec-ppt}, we now propose the \textbf{autocorrelated
multi-step conformal prediction} (AcMCP) method. Unlike extensions of
existing conformal prediction methods that treat multi-step forecasting
as independent events (see Section~\ref{sec-ext}), the AcMCP method
integrates the autocorrelations inherent in multi-step forecast errors,
thereby making the output multi-step prediction intervals more logically
structured.

The AcMCP method updates the quantile estimate \(q_t\) in an online
setting to achieve the goal of long-run coverage. Specifically, the
iteration of the \(h\)-step-ahead quantile estimate is given by

\begin{equation}\phantomsection\label{eq-acmcp}{
q_{t+h|t}=q_{t+h-1|t-1}+\eta \left(\mathrm{err}_{t|t-h}-\alpha\right)+r_t\left(\sum_{i=h+1}^t \left(\mathrm{err}_{i|i-h}-\alpha\right)\right)+\tilde{e}_{t+h|t},
}\end{equation}

for \(h\in[H]\). Obviously, the AcMCP method can be viewed as a further
extension of the PID method. Nevertheless, AcMCP diverges from PID with
several innovations and differences.

First, we are no longer confined to predicting just one step forward.
Instead, we can make multi-step forecasting, constructing
distribution-free prediction intervals for steps \(t+1,\ldots,t+h\)
based on available information up to time \(t\). This is highly
important in the field of time series forecasting.

Additionally, in AcMCP, \(\tilde{e}_{t+h|t}\) is a forecast combination
of two simple models: one being an MA\((h-1)\) model trained on
\(h\)-step-ahead forecast errors available up to and including time
\(t\) (i.e.~\(e_{1+h|1}, \ldots, e_{t|t-h}\)), and the other an
AR\((h-1)\) model (with respect to \(h\) instead of \(t\)) trained by
regressing \(e_{t+h|t}\) on forecast errors from past steps
(i.e.~\(e_{t+h-1|t}, \ldots, e_{t+1|t}\)). Thus, we perform multi-step
conformal prediction recursively, contrasting with the independent
approach employed in MPID. Moreover, the inclusion of
\(\tilde{e}_{t+h|t}\) is not intended to forecast the nonconformity
scores (i.e., forecast errors in this paper), but rather to incorporate
autocorrelations present in multi-step forecast errors within the
resulting multi-step prediction intervals. As previously explained, in
the context of time series forecasting, we typically assume the use of a
good base forecasting model, making it unnecessary to train an
additional model to predict forecast errors in order to capture leftover
information. If the forecasts are not optimal, the base forecasting
model should be improved to enhance its performance.

\subsubsection{Coverage guarantees}\label{coverage-guarantees}

\begin{proposition}[]\protect\hypertarget{prp-cov_rt}{}\label{prp-cov_rt}

Let \(\{s_{t+h|t}\}_{t\in\mathbb{N}}\) be any sequence of numbers in
\([-b, b]\) for any \(h\in[H]\), where \(b>0\), and may be infinite.
Assume that \(r_t\) is a saturation function obeying
Equation~\ref{eq-saturation_h}, for an admissible function \(g\). Then
the iteration
\(q_{t+h|t}=r_t\left(\sum_{i=h+1}^t\left(\mathrm{err}_{i|i-h}-\alpha\right)\right)\)
satisfies

\begin{equation}\phantomsection\label{eq-cov_rt}{
\left|\frac{1}{T-h}\sum_{t=h+1}^{T}\left(\mathrm{err}_{t|t-h}-\alpha\right)\right| \leq \frac{c \cdot g(T-h) + h}{T-h},
}\end{equation}

for any \(T \geq h+1\), where \(c>0\) is the constant in
Equation~\ref{eq-saturation_h}.

In particular, this means the prediction intervals obtained by the
iteration yield long-run coverage, i.e.,
\(\lim _{T \rightarrow \infty} \frac{1}{T-h} \sum_{t=h+1}^T \mathrm{err}_{t|t-h} = \alpha\).

\end{proposition}

\begin{proof}
Let \(E_T=\sum_{t=h+1}^{T}\left(\mathrm{err}_{t|t-h}-\alpha\right)\).
The inequality given by Equation~\ref{eq-cov_rt} can be expressed as
\(\left|E_T\right| \leq c \cdot g(T-h) + h\). We will prove one side of
the absolute inequality, specifically \(E_T \leq c \cdot g(T-h) + h\),
with the other side following analogously. We proceed with the proof
using induction.

For \(T=h+1,\ldots,2h\),
\(E_T = \sum_{t=h+1}^{T}(\mathrm{err}_{t|t-h}-\alpha) \leq (T-h)-(T-h)\alpha \leq T-h \leq h \leq cg(T-h) + h\)
as \(c>0\), \(h\geq 1\), \(g\) is nonnegative, and
\(\mathrm{err}_{t|t-h} \leq 1\). Thus, Equation~\ref{eq-cov_rt} holds
for \(T=h+1,\ldots,2h\).

Now, assuming Equation~\ref{eq-cov_rt} is true up to \(T\). We partition
the argument into \(h+1\) cases:

\[
\begin{cases}
cg(T-h)+h-1 < E_T \leq cg(T-h)+h, & \ldots \text { case (1) } \\
cg(T-h)+h-2 < E_T \leq cg(T-h)+h-1, & \ldots \text { case (2) } \\
\qquad \cdots \\
cg(T-h) < E_T \leq cg(T-h)+1, & \ldots \text { case (h) } \\
E_T \leq cg(T-h). & \ldots \text { case (h+1) }
\end{cases}
\]

In case (1), we observe that \(E_T > cg(T-h)+h-1 > cg(T-h)\), implying
\(q_{T+h|T} = r_t(E_{T}) \geq b\) according to
Equation~\ref{eq-saturation_h}. Thus, \(s_{T+h|T} \leq q_{T+h|T}\) and
\(\mathrm{err}_{T+h|T} = 0\). Furthermore, we have
\(E_{T-1} = E_T - (\mathrm{err}_{T|T-h} - \alpha) > cg(T-h)+h-2 > cg(T-h-1)\)
as \(g\) is nondecreasing. This implies
\(q_{T+h-1|T-1} = r_t(E_{T-1}) \geq b\), hence
\(s_{T+h-1|T-1} \leq q_{T+h-1|T-1}\) and
\(\mathrm{err}_{T+h-1|T-1} = 0\). Similarly,
\(E_{T-2} = E_{T-1} - (\mathrm{err}_{T-1|T-h-1} - \alpha) > cg(T-h)+h-3 > cg(T-h-2)\),
thus \(\mathrm{err}_{T+h-2|T-2} = 0\). This process iterates, leading to
\(\mathrm{err}_{T+h|T} = \mathrm{err}_{T+h-1|T-1} = \cdots = \mathrm{err}_{T+1|T-h+1} = 0\).
Consequently,

\[
E_{T+h} = E_T+\sum_{t=T+1}^{T+h}(\mathrm{err}_{t|t-h}-\alpha) \leq cg(T-h)+h-h\alpha \leq cg(T)+h,
\]

which is the desired result at \(T+h\).

In case (2), we observe that \(E_T > cg(T-h)+h-2 > cg(T-h)\), thus
\(s_{T+h|T} \leq q_{T+h|T}\) and \(\mathrm{err}_{T+h|T} = 0\). Moving
forward, we have
\(\mathrm{err}_{T+h|T} = \mathrm{err}_{T+h-1|T-1} = \cdots = \mathrm{err}_{T+2|T-h+2} = 0\).
Along with \(\mathrm{err}_{T+1|T-h+1} \leq 1\), this means that

\[
E_{T+h} = E_T+\sum_{t=T+1}^{T+h}(\mathrm{err}_{t|t-h}-\alpha) \leq cg(T-h)+h-1+1-h\alpha \leq cg(T)+h,
\]

which again gives the desired result at \(T+h\).

Similarly, in cases (3)-(h), we can always get the desired result at
\(T+h\).

In case (h+1), noticing \(E_T \leq cg(T-h)\), and simply using
\(\mathrm{err}_{T+h-i|T-i} \leq 1\) for \(i=0,\ldots,h-1\), we have

\[
E_{T+h} = E_T+\sum_{t=T+1}^{T+h}(\mathrm{err}_{t|t-h}-\alpha) \leq cg(T-h)+h-h\alpha \leq cg(T)+h.
\]

Therefore, we can deduce the desired outcome at any \(T \geq h+1\). This
completes the proof for the first part of Proposition~\ref{prp-cov_rt}.

Regarding the second part, \(g(t-h)/(t-h) \rightarrow 0\) as
\(t \rightarrow \infty\) due to the sublinearity of the admissible
function \(g\). Hence the second part holds trivially.
\end{proof}

The quantile iteration
\(q_{t+h|t}=q_{t+h-1|t-1}+\eta \left(\mathrm{err}_{t|t-h}-\alpha\right)\)
can be seen as a particular instance of the iteration outlined in
Proposition~\ref{prp-cov_rt} if we set \(q_{2h|h}=0\) without losing
generality. Thus, its coverage bounds can be easily derived as a result
of Proposition~\ref{prp-cov_rt}.

\begin{proposition}[]\protect\hypertarget{prp-cov_qt}{}\label{prp-cov_qt}

Let \(\{s_{t+h|t}\}_{t\in\mathbb{N}}\) be any sequence of numbers in
\([-b, b]\) for any \(h\in[H]\), where \(b>0\), and may be infinite.
Then the iteration
\(q_{t+h|t}=q_{t+h-1|t-1}+\eta \left(\mathrm{err}_{t|t-h}-\alpha\right)\)
satisfies

\[
\left|\frac{1}{T-h}\sum_{t=h+1}^{T}\left(\mathrm{err}_{t|t-h}-\alpha\right)\right| \leq \frac{b + \eta h}{\eta\left(T-h\right)},
\]

for any learning rate \(\eta > 0\) and \(T \geq h+1\).

In particular, this means the prediction intervals obtained by the
iteration yield long-run coverage, i.e.,
\(\lim _{T \rightarrow \infty} \frac{1}{T-h} \sum_{t=h+1}^T \mathrm{err}_{t|t-h} = \alpha\).

\end{proposition}

\begin{proof}
We set \(q_{2h|h}=0\) without losing generality, the iteration
\(q_{t+h|t}=q_{t+h-1|t-1}+\eta \left(\mathrm{err}_{t|t-h}-\alpha\right)\)
simplifies to
\(q_{t+h|t}=\eta \sum_{i=h+1}^{t}\left(\mathrm{err}_{i|i-h}-\alpha\right)\).
Let \(r_t(x) = \eta x\) and the admissible function \(g(t-h) = b\),
Equation~\ref{eq-saturation_h} holds for \(c=\frac{1}{\eta}\). Then
Proposition~\ref{prp-cov_rt} applies and we can easily derive the
desired result.
\end{proof}

More importantly, Proposition~\ref{prp-cov_rt} is also adequate for
establishing the coverage guarantee of the proposed AcMCP method given
by Equation~\ref{eq-acmcp}. Following the idea of
\textcite{angelopoulos2024}, we first reformulate
Equation~\ref{eq-acmcp} as

\begin{equation}\phantomsection\label{eq-acmcp_1}{
q_{t+h|t}=\hat{q}_{t+h|t}+r_t\left(\sum_{i=h+1}^t \left(\mathrm{err}_{i|i-h}-\alpha\right)\right),
}\end{equation}

where \(\hat{q}_{t+h|t}\) can be any function of the past observations
\(\{(\bm{x}_i, y_i)\}_{1 \leq i \leq t}\) and quantile estimates
\(q_{i+h|i}\) for \(i \leq t-1\). Taking
\(\hat{q}_{t+h|t}=q_{t+h-1|t-1}+\eta \left(\mathrm{err}_{t|t-h}-\alpha\right)+\tilde{e}_{t+h|t}\)
will recover Equation~\ref{eq-acmcp}. We can consider
\(\hat{q}_{t+h|t}\) as the forecast of the quantile \(q_{t+h|t}\) based
on available historical data. We then present the coverage guarantee for
AcMCP given by Equation~\ref{eq-acmcp_1}.

\begin{proposition}[]\protect\hypertarget{prp-cov_acmcp}{}\label{prp-cov_acmcp}

Let \(\{\hat{q}_{t+h|t}\}_{t\in\mathbb{N}}\) be any sequence of numbers
in \([-\frac{b}{2}, \frac{b}{2}]\), and
\(\{s_{t+h|t}\}_{t\in\mathbb{N}}\) be any sequence of numbers in
\([-\frac{b}{2},\frac{b}{2}]\), for any \(h\in[H]\), \(b>0\) and may be
infinite. Assume that \(r_t\) is a saturation function obeying
Equation~\ref{eq-saturation_h}, for an admissible function \(g\). Then
the prediction intervals obtained by the AcMCP iteration given by
Equation~\ref{eq-acmcp_1} yield long-run coverage, i.e.,
\(\lim _{T \rightarrow \infty} \frac{1}{T-h} \sum_{t=h+1}^T \mathrm{err}_{t|t-h} = \alpha\).

\end{proposition}

\begin{proof}
Let \(q_{t+h|t}^{*}=q_{t+h|t}-\hat{q}_{t+h|t}\), then
Equation~\ref{eq-acmcp_1} transforms into an update process
\(q_{t+h|t}^{*}=r_t\left(\sum_{i=h+1}^t \left(\mathrm{err}_{i|i-h}-\alpha\right)\right)\),
which is an update with respect to \(q_{t+h|t}^{*}\). Under this new
framework, the nonconformity score becomes
\(s_{t+h|t}^{*}=s_{t+h|t}-\hat{q}_{t+h|t}\), with values ranging in
\([-b,b]\), given the assumption that both \(s_{t+h|t}\) and
\(\hat{q}_{t+h|t}\) fall within \([-\frac{b}{2},\frac{b}{2}]\). Thus,
Proposition~\ref{prp-cov_qt} can be directly applied to establish the
long-run coverage achieved by the AcMCP method.
\end{proof}

\begin{itemize}
\item
  Remark on propositions.
\item
  Choosing the learning rate? Analysis of the parameter effect.
\end{itemize}

\section{Experiments}\label{experiments}

In this section, we examine the empirical performance of the previously
proposed multi-step conformal prediction methods using two simulated
data settings and two real data examples. Throughout the experiments, we
adhere to the following parameter settings: we focus on the target
coverage level \(1-\alpha=0.9\); for the MWCP method, we use \(b=0.99\)
as per \textcite{barber2023}; following \textcite{angelopoulos2024}, we
use a step size parameter of \(\gamma=0.005\) for the MACP method, a
Theta model as the scorecaster in the MPID method, and a learning rate
of \(\eta=0.01\hat{B}_t\) for quantile tracking in the MPID and AcMCP
methods, where
\(\hat{B}_t=\max\{s_{t-\Delta+1|t-\Delta-h+1},\cdots,s_{t|t-h}\}\) is
the highest score over a tailing window and the window length \(\Delta\)
is set to be same as the length of the calibration set; we consider a
clipped version of MACP by imputing infinite intervals with the largest
score seen so far.

\subsection{Simulated examples}\label{simulated-examples}

\subsubsection{Linear autoregressive
process}\label{linear-autoregressive-process}

We first consider a simulated stationary time series which is generated
from an AR\((2)\) process

\[
y_t = 0.8y_{t-1} - 0.5y_{t-2} + \epsilon_t,
\]

where \(\epsilon_t\) is white noise with error variance
\(\sigma^2 = 1\). After an appropriate burn-in period, we generate
\(N=5000\) data points. Under the sequential split and online learning
settings, we create training sets \(\mathcal{D}_{\text{tr}}\) and
calibration sets \(\mathcal{D}_{\text{cal}}\), each with a length of
\(500\). We use AR\((2)\) models to generate \(1\)- to \(3\)-step-ahead
point forecasts (i.e.~\(H=3\)) with the automated algorithm implemented
in the \textbf{forecast} R package \autocite{hyndman2024}. The goal is
to generate prediction intervals using various proposed conformal
prediction methods and evaluate whether they can achieve the nominal
long-run coverage for each separate forecast horizon.

Figure~\ref{fig-AR2_cov} presents the rolling coverage and interval
width of each method for each forecast horizon, with metrics computed
over a rolling window of size \(500\). In terms of coverage, we observe
that MPID and AcMCP achieve approximately the desired \(90\%\) coverage
level over the rolling windows, while other methods, including the AR
model, undergo much wider swings away from the desired level, showing
large troughs and peaks in coverage as time changes. Turning to the
prediction interval width, the trajectories of the rolling mean and
median of the interval widths for each method are largely consistent.
AcMCP constructs narrower prediction intervals than MPID, despite both
methods achieving similar coverage. Moreover, we see that AcMCP tends to
offer adaptive prediction intervals, and results in wider intervals
especially when competing methods undercover, which is to be expected.
In short, AcMCP intervals offer greater adaptivity and more precise
coverage compared to AR, MSCP, MWCP and MACP. However, MPID achieves
tight coverage but at the cost of constructing wider prediction
intervals. This is due to the fact that the inclusion of a second model
(scorecaster) is likely to introduce large variance into the generated
prediction intervals. The results can be further elucidated with
Figure~\ref{fig-AR2_box}, which presents boxplots of rolling coverage
and interval width for each method and each forecast horizon.

\begin{figure}

\centering{

\includegraphics{cpts_files/figure-pdf/fig-AR2_cov-1.pdf}

}

\caption{\label{fig-AR2_cov}AR(2) simulation results showing rolling
coverage, mean and median interval width for each forecast horizon. The
displayed curves are smoothed over a rolling window of size \(500\). The
black dashed line indicates the target level of \(1-\alpha=0.9\).}

\end{figure}%

\begin{figure}

\centering{

\includegraphics{cpts_files/figure-pdf/fig-AR2_box-1.pdf}

}

\caption{\label{fig-AR2_box}AR(2) simulation results showing boxplots of
the rolling coverage and interval width for each method across different
forecast horizons. The red dashed lines show the target coverage level,
while the blue dashed lines indicate the median interval width of the
AcMCP method.}

\end{figure}%

We note that the inclusion of the last term \(\tilde{e}_{t+h|t}\) in
AcMCP should only result in a slight difference compared to the version
without this term, which we henceforth refer to as MPI. This is because,
the inclusion of \(\tilde{e}_{t+h|t}\) aims to capture autocorrelations
inherent in multi-step forecast errors and focuses on the mean of
forecast errors, whereas the whole update of AcMCP operates on quantiles
of scores. To illustrate the subtle difference in their results and
explore their origins, we visualize their prediction interval over a
truncated period of length \(500\), as shown in
Figure~\ref{fig-AR2_timeplot}. We observe that AcMCP and MPI indeed
construct similar prediction intervals so their lower and upper bounds
mostly overlap with each other. The main differences may occur around
the time 1320 and during the period 1470-1500, where AcMCP tends to have
a fanning-out effect, increasing the interval width as the forecast
horizon increases, compared to MPI. Figure~\ref{fig-AR2_timeplot} also
presents the prediction interval bounds given by MACP. The prediction
intervals of both AcMCP and MACP can capture certain patterns in the
actual observations, and there is no consistent pattern indicating
dominance of one method over the other in terms of interval width.

\begin{figure}

\centering{

\includegraphics{cpts_files/figure-pdf/fig-AR2_timeplot-1.pdf}

}

\caption{\label{fig-AR2_timeplot}AR(2) simulation results showing the
prediction interval bounds for the MACP, MPI, and AcMCP methods over a
truncated period of length 500.}

\end{figure}%

\subsubsection{Nonlinear autoregressive
process}\label{nonlinear-autoregressive-process}

We then consider the case of a nonlinear data generation process, which
happens in many practical time series applications. We specify the true
data generation process as

\[
y_t = \sin(y_{t-1}) + 0.5\log(y_{t-2} + 1) + 0.1y_{t-1}x_{1,t} + 0.3x_{2,t} + \epsilon_{t},
\]

where \(x_{1,t}\) and \(x_{2,t}\) are uniformly distributed on
\([0,1]\), and \(\epsilon_{t}\) is white noise with error variance
\(\sigma^2 = 0.1\). Thus, the time series \(y_t\) nonlinearly depends on
its lagged values \(y_{t-1}\) and \(y_{t-2}\), as well as exogenous
variables \(x_{1,t}\) and \(x_{2,t}\).

After an appropriate burn-in period, we generate \(N=2000\) data points.
Under the sequential split and online learning settings, we create
training sets \(\mathcal{D}_{\text{tr}}\) and calibration sets
\(\mathcal{D}_{\text{cal}}\), each with a length of \(500\). Given the
nonlinear structure of the data generation process, we use feed-forward
neural networks with a single hidden layer and lagged inputs to generate
\(1\)- to \(3\)-step-ahead point forecasts (i.e.~\(H=3\)) with the
automated algorithm implemented in the \textbf{forecast} R package. Note
that it is not straightforward for neural networks to derive interval
forecasts, thus we do not include neural network models when presenting
the results.

Figure~\ref{fig-NL_cov} illustrates the rolling coverage and interval
width of each method, with calculations based on a rolling window of
size \(100\). We see that MPID and AcMCP are able to maintain minor
fluctionations around the target coverage of \(90\%\) across all time
indices, contrasting with MSCP, MWCP, and MACP, which struggle to
sustain the target coverage and display pronounced fluctuations over
time. Moreover, all conformal prediction methods, except for MSCP,
construct adaptive prediction intervals. They widen intervals in
response to undercoverage and narrow them when overcoverage occurs.
Notably, MPID and AcMCP demonstrate greater adaptability, displaying
higher variability in interval widths compared to competing methods in
order to uphold the desired coverage. Lastly, AcMCP intervals are
evidently narrower than MPID intervals for \(2\)-step-ahead forecasting
but wider for \(3\)-step-ahead forecasting. AcMCP intervals appear to be
more reasonable, as they tend to widen with increasing forecast
horizons.

\begin{figure}

\centering{

\includegraphics{cpts_files/figure-pdf/fig-NL_cov-1.pdf}

}

\caption{\label{fig-NL_cov}Nonlinear simulation results showing rolling
coverage, mean and median interval width for each forecast horizon. The
displayed curves are smoothed over a rolling window of size \(100\). The
black dashed line indicates the target level of \(1-\alpha=0.9\).}

\end{figure}%

The most basic requirement for a prediction interval is to cover the
actual value with the desired probability. Evaluating the efficiency of
forecast intervals is relevant only when they demonstrate valid
coverage. As discussed, MSCP, MWCP and MACP have coverage that deviate a
bit far from the desired coverage, thus we now focus our comparison on
MPID and AcMCP, the two methods that offer more accurate coverage.
Table~\ref{tbl-NL_winkler} shows the resulting coverage gap (the
difference between actual and nominal coverage), forecast interval
width, and Winkler score \autocite{winkler1972}, averaged over all test
sets, for the MPI, MPID, and AcMCP methods. The Winkler score is
proposed by \textcite{winkler1972} to evaluate a prediction interval,
and is defined as the length of the interval plus a penalty if the
observation is outside the interval. A smaller Winkler score indicates
better performance. The results indicate that, in this nonlinear data
generation process setting, there is essentially no difference among
these three methods in terms of overall coverage gap, average interval
width, and Winkler score.

\begin{table}

\caption{\label{tbl-NL_winkler}Nonlinear simulation results showing
coverage gap, interval width, and Winkler score, averaged over all test
sets.}

\centering{

\centering
\resizebox{\linewidth}{!}{
\fontsize{13}{15}\selectfont
\begin{tabular}{lrrrrrrrrr}
\toprule
\multicolumn{1}{c}{ } & \multicolumn{3}{c}{h=1} & \multicolumn{3}{c}{h=2} & \multicolumn{3}{c}{h=3} \\
\cmidrule(l{3pt}r{3pt}){2-4} \cmidrule(l{3pt}r{3pt}){5-7} \cmidrule(l{3pt}r{3pt}){8-10}
Methods & Coverage gap & Mean width & Winkler score & Coverage gap & Mean width & Winkler score & Coverage gap & Mean width & Winkler score\\
\midrule
MPI & 0 & 0.361 & 0.436 & 0.003 & 0.373 & 0.468 & 0.001 & 0.373 & 0.459\\
MPID & 0 & 0.362 & 0.436 & 0.003 & 0.388 & 0.476 & 0.001 & 0.367 & 0.460\\
AcMCP & 0 & 0.361 & 0.435 & 0.002 & 0.369 & 0.467 & 0.001 & 0.376 & 0.459\\
\bottomrule
\end{tabular}}

}

\end{table}%

We provide further insights into the performance of these conformal
prediction methods by presenting boxplots of the rolling coverage and
interval width for each method, as depicted in Figure~\ref{fig-NL_box}.
We see that coverage variability is higher for MSCP, MWCP and MACP than
for MPID and AcMCP, while MPID and AcMCP lead to a lower effective
interval size.

\begin{figure}

\centering{

\includegraphics{cpts_files/figure-pdf/fig-NL_box-1.pdf}

}

\caption{\label{fig-NL_box}Nonlinear simulation results showing boxplots
of the rolling coverage and interval width for each method across
different forecast horizons. The red dashed lines show the target
coverage level, while the blue dashed lines indicate the median interval
width of the AcMCP method.}

\end{figure}%

\subsection{Real data examples}\label{real-data-examples}

\subsubsection{Electricity demand data}\label{electricity-demand-data}

Now we examine empirical performance of the conformal prediction methods
using an electricity demand data set. The data set tracks daily
electricity demand (GW), daily maximum temperature (degrees Celsius),
and holiday information for the state of Victoria, Australia, spanning a
three-year period from 2012 to 2014. Figure~\ref{fig-elec_data} displays
the daily electricity demand during 2012-2014, along with temperatures.
The right panel shows a nonlinear relationship between electricity
demand and temperature, with demand increasing for low temperatures (due
to heating) and increasing for high temperatures (due to cooling).

\begin{figure}

\centering{

\includegraphics{cpts_files/figure-pdf/fig-elec_data-1.pdf}

}

\caption{\label{fig-elec_data}Daily electricity demand and corresponding
daily maximum temperatures in 2012--2014, Victoria, Australia.}

\end{figure}%

Our response variable is \(\texttt{Demand}\), and we use two covariates:
\(\texttt{Temperature}\), and \(\texttt{Workday}\) which is an indicator
variable for if the day was a working day or not. Following
\textcite{hyndman2021}, we will fit a dynamic regression model with a
piecewise linear function of temperature (containing a knot at \(18\)
degrees) to generate \(1\)- to \(7\)-step-ahead point forecasts
(i.e.~\(H=7\)). The error series in the regression is assumed to follow
an ARIMA model to contain autocorrelation. Under the sequential split
and online learning settings, we use two years of data as training sets
to fit dynamic regression models, and use \(100\) data points for
calibration sets.

We present the results in Figure~\ref{fig-elec_cov} and
Figure~\ref{fig-elec_box}, comparing the rolling coverage and interval
width of each method. These computations are based on a rolling window
of size \(100\). First, we observe that DR (dynamic regression)
consistently achieves a significantly higher coverage than the \(90\%\)
target coverage, resulting in much wider intervals than other methods
for \(h=1,2,3,4\). Secondly, MSCP, MWCP, and MACP fail to sustain the
target coverage and noticeably undercover after September 2014 for all
forecast horizons, thus leading to narrower intervals than others.
Thirdly, MPI, MPID, and AcMCP offer prediction intervals that are wider
than those of other conformal prediction methods, effectively mitigating
or avoiding the undercoverage issue observed after September 2014.
Additionally, we notice that MPID performs slightly worse than MPI and
AcMCP in terms of coverage for \(h=3,4,7\), despite leading to wider
intervals. Finally, MPI and AcMCP coverage display similar pattern, but
AcMCP is capable of constructing narrower intervals than MPI.

\begin{figure}

\centering{

\includegraphics{cpts_files/figure-pdf/fig-elec_cov-1.pdf}

}

\caption{\label{fig-elec_cov}Electricity demand data results showing
rolling coverage, mean and median interval width for each forecast
horizon. The dynamic regression model is abbreviated to DR. The
displayed curves are smoothed over a rolling window of size \(100\). The
black dashed line indicates the target level of \(1-\alpha=0.9\).}

\end{figure}%

\begin{figure}

\centering{

\includegraphics{cpts_files/figure-pdf/fig-elec_box-1.pdf}

}

\caption{\label{fig-elec_box}Electricity demand data results showing
boxplots of the rolling coverage and interval width for each method
across different forecast horizons. The dynamic regression model is
abbreviated to DR. The red dashed lines show the target coverage level,
while the blue dashed lines indicate the median interval width of the
AcMCP method.}

\end{figure}%

We present the forecast interval bounds for MACP, MPI, and AcMCP in
Figure~\ref{fig-elec_timeplot}. The plot shows that MACP intervals are
too narrow to adequately hug the true values, particularly from
September to November of 2014. In contrast, MPI and AcMCP perform better
by widening their intervals, resulting in narrower swings away from the
\(90\%\) target level. The differences between MPI and AcMCP prediction
intervals are most pronounced in the \(5\)-, \(6\)-, and
\(7\)-step-ahead forecast results. For \(5\)-step-ahead forecasting,
from May to July of 2014, the upper bounds of MPI intervals struggle to
capture the true values. AcMCP addresses this undercoverage by construct
larger upper bounds. In August and September, AcMCP constructs smaller
upper bounds than MPI while still capturing the true values effectively.
For \(6\)-step-ahead forecasting from May to July, AcMCP offers smaller
upper bounds than MPI, which provides upper bounds that are far away
from the truth. Similar reaction is observed for \(7\)-step-ahead
forecasting during August and September.

\begin{figure}

\centering{

\includegraphics{cpts_files/figure-pdf/fig-elec_timeplot-1.pdf}

}

\caption{\label{fig-elec_timeplot}Electricity demand data results
showing the forecast interval bounds for MACP, MPI, and AcMCP over the
whole test set.}

\end{figure}%

Figure~\ref{fig-elec_winkler} presents Winkler scores of each method
across different forecast horizons. DR and MPID consistently perform
worse in terms of Winkler scores. For \(1\)-, \(2\)-, and \(3\)-step
ahead forecasting, MPI and AcMCP deliver comparable or even better
results than MSCP, MWCP, and MACP, while still providing more precise
coverage. For \(4\)-, \(5\)-, \(6\)-, and \(7\)-step ahead forecasting,
MPI and AcMCP yield much larger Winkler scores than other conformal
prediction methods, whereas the competing methods suffer from severe
undercoverage issues. Thus, MPI and AcMCP are able to provide valid
prediction intervals. We also notice that AcMCP outperforms MPI for
\(h=6\) and \(h=7\), where both methods offer similar coverage.

\begin{figure}

\centering{

\includegraphics{cpts_files/figure-pdf/fig-elec_winkler-1.pdf}

}

\caption{\label{fig-elec_winkler}Electricity demand data results showing
Winkler scores of each method across different forecast horizons.}

\end{figure}%

\subsubsection{Eating out expenditure
data}\label{eating-out-expenditure-data}

Finally, we apply the conformal prediction methods to predict the eating
out expenditure (\$million) in Victoria, Australia. The data set
includes monthly expenditure on cafes, restaurants and takeaway food
services in Victoria from April 1982 up to December 2019, as shown in
Figure~\ref{fig-cafe_data}. The data shows an overall upward trend,
obvious anuual seasonal patterns, and variability proportional to the
data level.

\begin{figure}

\centering{

\includegraphics{cpts_files/figure-pdf/fig-cafe_data-1.pdf}

}

\caption{\label{fig-cafe_data}Monthly expenditure on cafes, restaurants
and takeaway food services in Victoria, Australia, from April 1982 to
December 2019.}

\end{figure}%

We consider three models: ARIMA with logarithmic transformation, ETS,
and STL-ETS \autocite{hyndman2021}, and then output their simple average
as final point forecasts. STL-ETS involves forecasting using the STL
decomposition method, applying ETS to forecast the seasonally adjusted
series. All three models can be automatically trained using the
\textbf{forecast} R package. Our goal is to forecast \(12\) months
ahead, i.e.~\(H=12\). Under the sequential split and online learning
settings, we use \(20\) years of data for training the models and \(5\)
years of data points for calibration sets. The whole test set only has a
length of \(152\) months. Therefore, we will not compute rolling
coverage and interval width in this experiment, but rather compute the
coverage and interval width averaged over the entire test set.

We summarize the average coverage, interval width, and Winkler score for
each method and each forecast horizon in Figure~\ref{fig-cafe_result}.
The results first show that MSCP, MWCP and MACP provide valid prediction
intervals for smaller forecast horizon but fail to achieve the desired
coverage for larger forecast horizons (\(h>5\)). Secondly, for
\(h \leq 5\), MPI and AcMCP can approximately achieve the desired
coverage and provide comparable mean interval widths and Winkler scores
with other methods, except for MPID. Thirdly, the coverage of MSCP, MWCP
and MACP declines gradually as the forecast horizon increases, while MPI
and AcMCP maintain coverage within a tighter range, albeit at the cost
of interval efficiency. Lastly, compared to MPI, AcMCP exhibits less
deviation from the desired coverage across most forecast horizons.

\begin{figure}

\centering{

\includegraphics{cpts_files/figure-pdf/fig-cafe_result-1.pdf}

}

\caption{\label{fig-cafe_result}Eating out expenditure data results
showing coverage gap, interval width, and Winkler scores averaged over
the entire test set for each forecast horizon. The black dashed line in
the top panel indicates no difference from the \(90\%\) target level.}

\end{figure}%

\section{Discussion}\label{discussion}

\begin{itemize}
\item
  Large forecast horizon
\item
  Learning rate determination
\item
  Variability introduced by the forecasts of predictors
\item
  Efficiency of prediction intervals
\end{itemize}

\section*{References}\label{references}
\addcontentsline{toc}{section}{References}

\printbibliography[heading=none]

\newpage
\appendix
% \pagenumbering{arabic}% resets `page` counter to 1
\setcounter{section}{0}
\renewcommand{\thesection}{Appendix \Alph{section}}
\renewcommand{\thesubsection}{\Alph{section}.\arabic{subsection}}
\renewcommand{\thefigure}{A\arabic{figure}}
\renewcommand{\thetable}{A\arabic{table}}
\setcounter{figure}{0}
\setcounter{table}{0}

\section{Weighted quantile estimation}\label{sec-wq}

\begin{itemize}
\tightlist
\item
  median-unbiased estimates
\end{itemize}

\section{Proof}\label{sec-proof}

\subsection{\texorpdfstring{Proof of
Proposition~\ref{prp-ma}}{Proof of Proposition~}}\label{sec-proof_ma}

\begin{proof}
Here, we give the proof of Proposition~\ref{prp-ma} based on
Proposition~\ref{prp-ar}, and the idea is motivated by
\textcite{sommer2023}.

For \(1\)-step-ahead optimal forecast, we have

\[
y_{t+1} = f_{t+1}(y_{(t-d+1):t},\bm{x}_{(t-k+1):(t+1)}) + \epsilon_{t+1},
\]

so \(e_{t+1|t}=\epsilon_{t+1}\).

Based on Proposition~\ref{prp-ar}, we have

\[
\begin{aligned}
e_{t+2|t} &= \epsilon_{t+2} + \phi_{1}^{(2)}e_{t+1|t} \\
e_{t+3|t} &= \epsilon_{t+3} + \phi_{1}^{(3)}e_{t+2|t} + \phi_{2}^{(3)}e_{t+1|t} \\
\cdots \\
e_{t+d|t} &= \epsilon_{t+d} + \phi_{1}^{(d)}e_{t+d-1|t} + \cdots + \phi_{d-1}^{(d)}e_{t+1|t} \\
e_{t+d+1|t} &= \epsilon_{t+d+1} + \phi_{1}^{(d+1)}e_{t+d|t} + \cdots + \phi_{d-1}^{(d+1)}e_{t+2|t} + \phi_{d}^{(d+1)}e_{t+1|t} \\
\cdots \\
e_{t+H|t} &= \epsilon_{t+H} + \phi_{1}^{(H)}e_{t+H-1|t} + \cdots + \phi_{d-1}^{(H)}e_{t+H-d+1|t} + \phi_{d}^{(H)}e_{t+H-d|t}, \quad H > d + 1. \\
\end{aligned}
\]

Substituting all equations above into the following equation, we can
obtain

\[
e_{t+h|t} = \epsilon_{t+h} + \sum_{i=1}^{h-1}\theta_{i}\epsilon_{t+h-i}, \text{ for each } h\in[H],
\]

where \(\theta_{i}\) is a complex function derived from the AR
coefficients of all AR\((j-1)\) models, for \(j = 1,2,\ldots,h-1\). So
we conclude that the \(h\)-step-ahead forecast error sequence
\(\{e_{t+h|t}\}\) follows an MA\((h-1)\) process.
\end{proof}




\end{document}
